\section{Introduction}

Pour le cours de théorie des graphes, nous (Alexandr Pavlov et Joachim Lardinois) avons du faire un programme capable de 3 choses.
Le programme est capable de lire un fichier contenant un graphe, et d'écrire dans un autre fichier un sous-arbre couvrant de poids minimal (MST) de ce graphe.
Ensuite, il doit pouvoir trouver un second MST différent du premier.
Finallement, en modifiant un graphe, notre programme doit mêtre à jour un MST de ce graphe, idéalement plus rapidement que en le reconstruisant completement.
